% Прямоходовой преобразователь (forward) — топология по Кобзеву, рис. 6.3.
% Ключ VT (MOSFET с индуцированным каналом); W_p — размагничивающая
% обмотка (нижний вывод с точкой — на плюсовом стыке с W_1, верхний —
% через VD3 на минусовую шину); VD1 выпрямительный, VD2 нулевой,
% дроссель L, C, R_н. Гальваническая развязка.
\documentclass[border=4pt]{standalone}
\usepackage{amsmath}
\usepackage{fontspec}
\setmainfont{Calibri}
\usepackage[european, RPvoltages]{circuitikz}
\usetikzlibrary{arrows.meta, calc}
\definecolor{ink}{RGB}{31,42,51}
\definecolor{grp}{RGB}{46,110,140}
\ctikzset{bipoles/length=1.0cm}
\ctikzset{bipoles/thickness=0.9}
\ctikzset{diodes/fill=ink}
\ctikzset{inductor=american}

\begin{document}
\begin{tikzpicture}[line width=1.0pt, color=ink, font=\small]
  % --- входные зажимы (открытые)
  \draw (0.6,2.6) node[ocirc]{} node[above left]{$+$} -- (2.2,2.6)
        node[circ]{};
  \node[anchor=east] at (0.4,1.3) {$U_{\text{вх}}$};
  % --- первичная обмотка W1 (нижняя) и размагничивающая Wp (верхняя),
  %     стык — плюсовой узел (точки одноимённых выводов у стыка)
  \draw (2.2,1.35) to[L, l=$w_1$] (2.2,2.55) -- (2.2,2.6);
  \draw (2.2,2.6) -- (2.2,2.65) to[L, l=$w_p$] (2.2,3.85);
  \fill (2.44,2.45) circle (1.6pt);
  \fill (2.44,3.6) circle (1.6pt);
  % --- ключ VT: корпус по центру между шиной и обмоткой
  \draw (2.2,0) node[circ]{} -- (2.2,0.45);
  \draw (2.2,0.45) to[Tnigfete, name=VT] (2.2,0.9) -- (2.2,1.1);
  \node[anchor=east] at (1.82,0.675) {VT};
  \draw[grp] (VT.G) -- ++(-0.5,0);
  \node[anchor=east, text=grp] at ($(VT.G)+(-0.55,0)$) {ШИМ};
  % --- нижняя шина первичной цепи
  \draw (0.6,0) node[ocirc]{} node[below left]{$-$} -- (2.2,0);
  % --- цепь размагничивания: минусовая шина -> VD3 (анод внизу,
  %     катод к верхнему выводу Wp) — по Кобзеву, рис. 6.3
  \draw (2.2,3.85) -- (2.2,4.3) -- (0.9,4.3) -- (0.9,2.8);
  \draw (0.9,1.9) to[D, name=VD3] (0.9,2.8);
  \draw (0.9,1.9) -- (0.9,0) node[circ]{};
  \node[anchor=west] at (1.15,2.35) {VD3};
  % --- магнитопровод (две вертикальные линии) + вторичная обмотка
  \draw[line width=0.8pt] (2.8,1.2) -- (2.8,4.0);
  \draw[line width=0.8pt] (2.92,1.2) -- (2.92,4.0);
  \draw (3.5,1.4) to[L, l_=$w_2$] (3.5,3.8);
  \fill (3.26,3.55) circle (1.6pt);
  % --- выпрямитель: VD1 от верхнего вывода w2
  \draw (3.5,3.8) -- (3.5,4.1) -- (3.9,4.1) to[D, name=VD1] (4.6,4.1)
        -- (4.9,4.1);
  \node[anchor=south] at (4.25,4.22) {VD1};
  % --- дроссель L в шине выхода
  \draw (4.9,4.1) to[L, l_=$L$] (6.0,4.1) node[circ]{} -- (8.6,4.1);
  % --- нулевой диод VD2: катод к узлу перед L (остриё вверх)
  \draw (4.75,4.1) -- (4.75,3.6);
  \draw (4.75,1.2) to[D, name=VD2] (4.75,3.6);
  \draw (4.75,1.2) -- (4.75,0.4) node[circ]{};
  \node[anchor=west] at (5.0,2.4) {VD2};
  % --- нижняя шина ВТОРИЧНОЙ цепи (развязка!)
  \draw (3.5,1.4) -- (3.5,0.4) -- (8.6,0.4);
  % --- конденсатор и нагрузка
  \draw (7.0,4.1) node[circ]{} to[C, l_=$C$] (7.0,0.4) node[circ]{};
  \draw (8.0,4.1) node[circ]{} to[R, l=$R_{\text{н}}$] (8.0,0.4)
        node[circ]{};
  % --- выходные зажимы
  \draw (8.6,4.1) node[ocirc]{} node[above right]{$+$};
  \draw (8.6,0.4) node[ocirc]{} node[below right]{$-$};
  \node[anchor=west] at (8.8,2.2) {$U_{\text{вых}}$};
\end{tikzpicture}
\end{document}
